\section{Execution Layer Position in AI Native Stack}
\label{sec:stack}

\subsection{Five-Layer Architecture}

The execution layer sits within a broader AI Native implementation stack:

\begin{figure}[t]
\centering
\begin{tikzpicture}[
  layer/.style={draw=black!35, rounded corners=3pt, minimum width=10.5cm, minimum height=1.1cm, align=center, font=\bfseries, inner sep=6pt},
  sublayer/.style={draw=black!30, rounded corners=2pt, minimum width=7.4cm, minimum height=0.95cm, align=center, font=\small, inner sep=4pt},
  flow/.style={-{Latex[length=2.2mm]}, thick, draw=black!55}
]
  \node[layer, fill=blue!10] (ui) at (0,0) {User Interface};
  \node[layer, fill=cyan!10, below=0.9cm of ui] (apps) {AI-Native Applications\\\normalfont(Agents, Workflows, Applications)};
  \node[layer, fill=teal!10, below=0.9cm of apps] (cog) {AI Cognition Layer\\\normalfont(LLMs, Reasoning, Planning)};
  \node[layer, fill=orange!10, below=1.1cm of cog, minimum height=3.1cm] (action) {AI Action Layer};
  \node[sublayer, fill=orange!20] (orch) at ([yshift=0.62cm]action.center) {Policy Orchestration Layer};
  \node[sublayer, fill=yellow!18] (runtime) at ([yshift=-0.72cm]action.center) {Execution Runtime};
  \node[layer, fill=purple!10, below=1.1cm of action] (eco) {Capability Ecosystem\\\normalfont(Models, APIs, Tools, Devices)};
  \node[layer, fill=gray!12, below=0.9cm of eco] (infra) {Infrastructure\\\normalfont(Cloud, Edge, Networks)};

  \draw[flow] (ui) -- (apps);
  \draw[flow] (apps) -- (cog);
  \draw[flow] (cog) -- (action);
  \draw[flow] (action) -- (eco);
  \draw[flow] (eco) -- (infra);
  \draw[flow] ([xshift=0.2cm]orch.south) -- ([xshift=0.2cm]runtime.north);
\end{tikzpicture}
\caption{Five-layer AI Native implementation stack.}
\label{fig:five-layer-stack}
\end{figure}

\subsection{Action Layer Decomposition}

A critical insight is that the Action Layer comprises two subcomponents:

\textbf{Orchestration Component (Contact Layer):}
\begin{itemize}[leftmargin=1.5em,nosep]
  \item Capability routing
  \item Provider selection
  \item Policy evaluation
  \item Environment selection
\end{itemize}

\textbf{Execution Component (Runtime):}
\begin{itemize}[leftmargin=1.5em,nosep]
  \item API invocation
  \item Tool execution
  \item Error handling
  \item Result normalization
\end{itemize}

This decomposition maintains clear separation of concerns. The orchestration layer makes \emph{decisions} (which provider, which environment), while the execution layer performs \emph{actions} (invoke API, return result).

When mapped to the logical model $S = (A, P, E)$, the Contact Layer corresponds to $P$, and the Execution Runtime corresponds to $E$.

\subsection{Defining Statement}

\begin{quote}
\emph{The execution layer is the action interface between AI cognition and the real-world capability ecosystem.}
\end{quote}

This definition captures the essential role of the layer: mediating between what agents decide to do and how those decisions are realized through capability execution.
